\documentclass[11pt]{article}

\usepackage[a4paper,margin=1in]{geometry}
\usepackage{amsmath,amssymb,amsthm,mathtools}
\usepackage{booktabs,array,longtable}
\usepackage{enumitem}
\usepackage{microtype}
\usepackage{xcolor}
\usepackage{hyperref}
\usepackage[nameinlink,capitalise,noabbrev]{cleveref}

\hypersetup{
  colorlinks=true,
  linkcolor=blue!55!black,
  citecolor=blue!55!black,
  urlcolor=blue!55!black
}

\newtheorem{theorem}{Theorem}[section]
\newtheorem{proposition}[theorem]{Proposition}
\newtheorem{lemma}[theorem]{Lemma}
\newtheorem{corollary}[theorem]{Corollary}
\newtheorem{definition}[theorem]{Definition}
\newtheorem{remark}[theorem]{Remark}

% The environments above share the theorem counter, so cleveref would name
% them all "Theorem"; these aliases restore the correct names.
\usepackage{etoolbox}
\AtBeginEnvironment{proposition}{\crefalias{theorem}{proposition}}
\AtBeginEnvironment{lemma}{\crefalias{theorem}{lemma}}
\AtBeginEnvironment{corollary}{\crefalias{theorem}{corollary}}
\AtBeginEnvironment{definition}{\crefalias{theorem}{definition}}
\AtBeginEnvironment{remark}{\crefalias{theorem}{remark}}

\numberwithin{equation}{section}

\newcommand{\one}{\mathbf 1}
\newcommand{\R}{\mathbb R}
\newcommand{\ip}[2]{\left\langle #1,#2\right\rangle}
\newcommand{\norm}[1]{\left\lVert #1\right\rVert}
\newcommand{\Tr}{\operatorname{Tr}}
\newcommand{\ord}{\operatorname{ord}}
\newcommand{\pos}[1]{\left(#1\right)_{+}}
\newcommand{\dd}{\,\mathrm d}
\newcommand{\E}{\mathbb E}
\newcommand{\GammaLaw}{\operatorname{Gamma}}
\newcommand{\BetaLaw}{\operatorname{Beta}}
\newcommand{\supp}{\operatorname{supp}}

\newcommand*{\tk}[1]{\textcolor{blue}{\textbf{[TK: #1]}}}
\newcommand*{\jl}[1]{\textcolor{red!70!black}{\textbf{[JL: #1]}}}
\newcommand*{\jb}[1]{\textcolor{green!50!black}{\textbf{[JB: #1]}}}
\newcommand*{\si}[1]{\textcolor{orange!80!black}{\textbf{[SI: #1]}}}
\newcommand*{\jk}[1]{\textcolor{purple}{\textbf{[JK: #1]}}}
\newcommand*{\hy}[1]{\textcolor{teal}{\textbf{[HY: #1]}}}

\newcommand{\dens}[2]{t(#1,#2)}
\newcommand{\tC}[2]{t(C_{#1},#2)}
\newcommand{\tP}[2]{t(P_{#1},#2)}
\newcommand{\edens}[1]{t(K_{2},#1)}

\newcommand{\goodman}[2]{#1^{#2}-#1(1-#1)^{#2-1}}
\newcommand{\goal}{\goodman{p}{m}}

\newcommand{\Tker}[1]{T_{#1}}
\newcommand{\dW}{d_W}

\newcommand{\coeff}[1]{[z^{#1}]}
\newcommand{\nld}[1]{-\log\det(I-z#1)}

\newcommand{\resm}{(I-zA)^{-1}}
\newcommand{\resp}{(I+zA)^{-1}}
\newcommand{\gresm}{\ip g{\resm g}}
\newcommand{\gresp}{\ip g{\resp g}}

\newcommand{\specmu}{\sum_i\norm{P_ig}^2\,\delta_{\lambda_i}}
\newcommand{\dmu}{\dd\mu(\lambda)}
\newcommand{\intmu}[1]{\int #1\dmu}
\newcommand{\gA}[1]{\ip g{A^{#1}g}}

\newcommand{\intsq}{\int_{[0,1]^2}}
\newcommand{\intcube}[1]{\int_{[0,1]^{#1}}}
\newcommand{\ddxy}{\dd x\dd y}

\newcommand{\rep}[2]{#1^{\times #2}}

\newcommand{\blockmat}[4]{\begin{pmatrix}#1&#2\\#3&#4\end{pmatrix}}
\newcommand{\sblockmat}[4]{\left(\begin{smallmatrix}#1&#2\\#3&#4\end{smallmatrix}\right)}

\title{On the Goodman style lower bounds of the odd cycles}
\author{%
  Taeyoung Kim\thanks{School of Computing, KAIST, Daejeon, South Korea.}
  \and Sunghyuk Im\thanks{Center for AI and Natural Science, KIAS, Seoul, South Korea.}
  \and Jineon Baek\thanks{June E Huh Center for Mathematical Challenges, KIAS, Seoul, South Korea.}
  \and Joonkyung Lee\thanks{Department of Mathematics, Yonsei University, Seoul, South Korea.}
  \and Jangsoo Kim\thanks{Department of Mathematics, Sungkyunkwan University, Suwon, South Korea.}
  \and Hongseok Yang\thanks{School of Computing, KAIST, Daejeon, South Korea, and School of Computational Sciences, KIAS, Seoul, South Korea.}%
}
\date{\today}

\begin{document}
% Symbolic marks (*, †, ‡, §, ¶, ‖) for the affiliation footnotes on the
% title page; regular numbered footnotes resume after \maketitle.
\renewcommand{\thefootnote}{\fnsymbol{footnote}}
\maketitle
\renewcommand{\thefootnote}{\arabic{footnote}}
\setcounter{footnote}{0}

\begin{abstract}
  Recently, Bennett et al. proved that the density of the odd cycle of length 5 is minimised by the balanced complete $k$-partite graph at the fixed edge density $p = 1 - 1 / k$, using flag algebras.
  Motivated by this result, we extend it to all odd cycles.
  Specifically, we prove that $f_m(p) = p^{m} - p(1-p)^{m-1}$ is a lower bound on the density of the odd cycle of length $m$ at every edge density $p \ge 2/3$, which is tight at $p = 1 - 1 / k$ for all $k \ge 2$.
  As a companion result, we also prove a stronger bound on the extended commonality of the odd cycles, which strengthens the result of Kim and Lee.
\end{abstract}

\tableofcontents

\section{Introduction}\label{sec:intro}

A \emph{graphon} is a symmetric measurable function $W\colon[0,1]^2\to[0,1]$.
For a finite graph $F$ with vertex set $\{1,\ldots,v\}$ and edge set $E(F)$, the \emph{homomorphism density} of $F$ in $W$ is
\[
  \dens FW=\intcube{v}\prod_{ij\in E(F)}W(x_i,x_j)\dd x_1\cdots\dd x_v,
\]
and $p=\edens W=\intsq W(x,y)\ddxy$ is the \emph{edge density} of $W$.

\subsection{The minimum density of an odd cycle}

How small can the density of a fixed graph be at a given edge density?
For the triangle, Goodman~\cite{goodman1959sets} proved that every graphon of edge density $p$ satisfies
\begin{equation}\label{eq:intro-goodman}
  \tC3W\ge2p^2-p=\goodman p3.
\end{equation}
For an integer $k\ge2$, equality holds at $p=1-1/k$ for the \emph{balanced complete $k$-partite graphon}
\[
  W_k(x,y)=
  \begin{cases}
    1&\text{if }\lfloor kx\rfloor\ne\lfloor ky\rfloor,\\
    0&\text{otherwise},
  \end{cases}
\]
whose edge density is $1-1/k$.
At the other edge densities the bound \eqref{eq:intro-goodman} is not attained, and the exact minimum---the subject of the Erd\H os--Rademacher problem~\cite{erdos1962number}---was determined only much later: Razborov~\cite{razborov2008minimal} found it for the triangle with his flag algebra method, Nikiforov~\cite{nikiforov2010number} for $K_4$, and Reiher~\cite{reiher2016clique} for every complete graph, confirming the extremal structure conjectured by Lov\'asz and Simonovits~\cite{lovasz1983number}.

Bipartite graphons contain no odd cycle and exist at every edge density $p\le1/2$, so for odd $m$ the minimum density of the cycle $C_m$ with $m$ vertices is $0$ at every $p\le1/2$, and the problem is nontrivial only for $p>1/2$.
There, the natural candidate for the minimizer at the density $p=1-1/k$ is again $W_k$.
Cycle densities are traces of kernel operators, $\tC mW=\Tr(\Tker W^m)=\sum_i\lambda_i^m$ for $m\ge3$, where $\Tker W$ is the integral operator on $L^2[0,1]$ with kernel $W$ and $\lambda_i$ are its eigenvalues~\cite{lovasz2012large}; the nonzero eigenvalues of $\Tker{W_k}$ are $p$, on the constant functions, and $-1/k$ with multiplicity $k-1$, so for odd $m$
\begin{equation}\label{eq:intro-kpartite-value}
  \tC m{W_k}=p^m+(k-1)\Bigl(-\frac1k\Bigr)^m=\goal.
\end{equation}
Whether this value is the minimum was known in two cases.
For $m=3$ it is Goodman's bound \eqref{eq:intro-goodman}.
For $m=5$, Bennett, Dudek, Lidick\'y and Pikhurko~\cite{bennett2020minimizing} proved that, for every integer $k\ge3$, a graph of edge density $1-1/k$ contains asymptotically at least as many copies of $C_5$ as the balanced complete $k$-partite graph; by the correspondence between graphs and graphons~\cite{lovasz2006limits,lovasz2012large}, $W_k$ minimizes the $C_5$-density among graphons of edge density $1-1/k$.
Their proof is a flag algebra computation together with a stability analysis.
Their theorem is stated only at the densities $p=1-1/k$, but interpolating their flag algebra certificate over real $k\in[3,\infty)$ yields the lower bound $\tC5W\ge\goodman p5$ at every edge density $p\ge2/3$.

Our main result is the corresponding lower bound for every odd cycle, at every edge density $p\ge2/3$.

\begin{theorem}[Odd-cycle Goodman bound for $p\ge2/3$]\label{thm:main}
  Let $m\ge3$ be odd, and let $W$ be a graphon of edge density $p\ge2/3$.
  Then
  \begin{equation}\label{eq:intro-main}
    \tC mW\ge\goal.
  \end{equation}
\end{theorem}

For $m=3$, \Cref{thm:main} is contained in Goodman's bound \eqref{eq:intro-goodman}, which holds at every edge density; the content of the theorem is the case $m\ge5$.
For $p\le1/2$ the right side of \eqref{eq:intro-main} is nonpositive, since $\goal=p(p^{m-1}-(1-p)^{m-1})$, so the inequality holds at every edge density $p\le1/2$ as well.
The remaining range $1/2<p<2/3$ needs different methods and is not treated in this paper; we discuss it in \Cref{sec:conclusion}.

By the value \eqref{eq:intro-kpartite-value}, the bound of \Cref{thm:main} is attained at the densities $p=1-1/k$.

\begin{corollary}\label{cor:kpartite}
  For every integer $k\ge2$ and every odd $m\ge3$, the balanced complete $k$-partite graphon $W_k$ minimizes the $C_m$-density among graphons of edge density $1-1/k$.
\end{corollary}

For $k\ge3$ the corollary follows from \Cref{thm:main} and the value \eqref{eq:intro-kpartite-value}; for $k=2$ it holds because $W_2$ is bipartite, so that $\tC m{W_2}=0$, while $\tC mW\ge0$ for every graphon.
The corollary extends the theorem of Bennett, Dudek, Lidick\'y and Pikhurko from $m=5$ to every odd cycle length.

\subsection{Commonality of odd cycles}

A graph $H$ with $e(H)$ edges is \emph{common} if
\[
  \dens HW+\dens H{1-W}\ge2^{1-e(H)}
\]
for every graphon $W$; equivalently, the number of monochromatic copies of $H$ in a $2$-colouring of the edges of a large complete graph is asymptotically minimized by the uniformly random colouring.
Goodman's theorem implies that the triangle is common.
Erd\H os~\cite{erdos1962number} conjectured that every complete graph is common, and Burr and Rosta~\cite{burr1980ramsey} that every graph is; Thomason~\cite{thomason1989disproof} disproved both conjectures, already for $K_4$, and Jagger, \v S\v tov\'i\v cek and Thomason~\cite{jagger1996multiplicities} proved that no graph containing $K_4$ is common.
Sidorenko~\cite{sidorenko1989cycles} proved that every cycle is common.
Kim and Lee~\cite{kim2024extended} proved the following \emph{extended commonality} of paths and cycles: for every path or cycle $H$ and every graphon $W$ of edge density $p$,
\begin{equation}\label{eq:intro-extended-commonality}
  \dens HW+\dens H{1-W}\ge p^{e(H)}+(1-p)^{e(H)},
\end{equation}
answering a question of Sidorenko~\cite{sidorenko1989cycles} and proving a conjecture of Behague, Morrison and Noel~\cite{behague2025common} in a strengthened form.
Since $p^{e}+(1-p)^{e}\ge2^{1-e}$, the extended inequality \eqref{eq:intro-extended-commonality} implies commonality.

For odd cycles at edge density $p\ge2/3$, we prove \eqref{eq:intro-extended-commonality} with an explicit nonnegative quantity added to its right side.

\begin{theorem}[Strengthened extended commonality]\label{thm:commonality}
  Let $m\ge3$ be odd, let $W$ be a graphon of edge density $p\ge2/3$, and write $U=1-W$ and $q=1-p$.
  Then, with $P_{m-1}$ the path with $m-1$ edges,
  \begin{equation}\label{eq:intro-commonality}
    \tC mW+\tC mU\ge p^m+q^m+\tP{m-1}U-q^{m-1}.
  \end{equation}
\end{theorem}

% NOTE: the path bound t(P_{m-1},U) >= q^{m-1} is the Mulholland--Smith inequality;
% it is also attributed to Blakley--Roy (1965); add that reference if preferred.
The added quantity is nonnegative: $\tP{m-1}U\ge q^{m-1}$ by the path inequality of Mulholland and Smith~\cite{mulholland1959inequality}.
Hence \eqref{eq:intro-commonality} strengthens \eqref{eq:intro-extended-commonality} for odd cycles whenever $p\ge2/3$; exchanging the roles of $W$ and $U$ gives the corresponding statement for $p\le1/3$.
Equality holds in \eqref{eq:intro-commonality} for the constant graphon $W\equiv p$ and, when $p=1-1/k$ with $k\ge3$, for $W_k$.

\Cref{thm:main} is a consequence of \Cref{thm:commonality}.
Since $0\le U\le1$, deleting one edge of the $m$-cycle gives $\tC mU\le\tP{m-1}U$, and subtracting this bound from \eqref{eq:intro-commonality} gives, with $p+q=1$,
\[
  \tC mW-\bigl(\goal\bigr)\ge\tP{m-1}U-\tC mU\ge0.
\]

% The overview below is commented out: the proof is going to be revised into a
% more combinatorially interpretable form, and this outline describes the old
% analytic route. Rewrite it once the revised proof structure is fixed.
\iffalse
\subsection{Overview of the proof}\label{subsec:intro-overview}

Every quantity in the proof is built from the complement kernel.
Let $U=1-W$, $q=1-p$, and $T=\Tker U$.
Splitting $T\one$ into its constant and mean-zero parts gives
\[
  T\one=q\one+g,
  \qquad
  g=p-\dW,
  \quad
  \dW(x)=\int_0^1W(x,y)\dd y,
\]
so $g$ is the deviation of the degree function of $W$ from its mean.
We restrict $T$ to the mean-zero subspace $\one^\perp$ and write $A$ for the resulting self-adjoint operator, in the sense that $\ip f{Ah}=\ip f{Th}$ for all $f,h\perp\one$, and we set $s_j=\ip g{A^jg}$.
Two elementary facts are used throughout.
First, $\Tker Uf=\bigl(\int_0^1f\bigr)\one-\Tker Wf$, so the operator obtained from $W$ in the same way is $-A$: the two kernels share one operator up to sign.
Second, since $0\le U\le1$, the spectrum of $A$ lies in $[-1/2,1/2]$.

After a reduction to step graphons, which fixes the edge density and makes every operator below finite-rank, cycle densities become traces, $\tC jK=\Tr(\Tker K^j)$ for $j\ge2$, and a determinant generates them:
\[
  \nld{\Tker K}=\sum_{j\ge1}\frac{\Tr(\Tker K^j)}{j}z^j.
\]
The Schur complement of $I-zT$ against the constant function factorizes the determinant into a scalar factor and a mean-zero factor,
\[
  \det(I-z\Tker U)=\bigl(1-qz-z^2R(z)\bigr)\det(I-zA),
  \qquad
  R(z)=\sum_{j\ge0}s_jz^j,
\]
and the same computation for $W$ gives $\det(I-z\Tker W)=\bigl(1-pz-z^2R(-z)\bigr)\det(I+zA)$.
Taking logarithms and extracting the coefficient of $z^m$ turns the two factorizations into two identities for the cycle densities, whose trace parts are $\Tr(A^m)$ and $\Tr((-A)^m)=-\Tr(A^m)$.
Adding the identities cancels the traces and leaves
\begin{equation}\label{eq:intro-two-sided}
  \tC mW+\tC mU=p^m+q^m+S_m,
\end{equation}
where $S_m$ is an explicit polynomial in $p$, $q$, and the moments $s_0,\ldots,s_{m-2}$.
The odd trace $\Tr(A^m)$, which neither kernel controls separately, does not appear in the identity \eqref{eq:intro-two-sided}.
By \eqref{eq:intro-two-sided}, the extended commonality \eqref{eq:intro-extended-commonality} for odd cycles is exactly the inequality $S_m\ge0$.
The path densities obey the same algebra: the scalar factor above generates them, $\sum_{j\ge0}\tP jUz^j=\bigl(1-qz-z^2R(z)\bigr)^{-1}$.
Therefore \Cref{thm:commonality} reduces to the single inequality
\[
  \Phi_m:=q^{m-1}+S_m-\tP{m-1}U\ge0.
\]

Positivity is proved spectrally.
Let $\lambda_i$ be the eigenvalues of $A$, let $P_i$ be the orthogonal projections onto their eigenspaces, and let $\mu=\specmu$, the spectral measure of $A$ associated with $g$; then $s_j=\intmu{\lambda^j}$, and $\mu$ is supported in $[-1/2,1/2]$.
Sorting the terms of $\Phi_m$ by the number $r$ of factors $s_j$ they contain gives, with $n=m-2r$,
\[
  \Phi_m=\sum_{r=1}^{(m-1)/2}\int\!\cdots\!\int
  P_{m,r}(q;\lambda_1,\ldots,\lambda_r)
  \prod_{i=1}^r\dd\mu(\lambda_i),
\]
\[
  P_{m,r}=\frac mr\Bigl[h_n(\rep pr,-\lambda_1,\ldots,-\lambda_r)
  +h_n(\rep qr,\lambda_1,\ldots,\lambda_r)\Bigr]
  -h_{n-1}(\rep q{r+1},\lambda_1,\ldots,\lambda_r),
\]
where $h_d$ is the complete homogeneous symmetric polynomial of degree $d$ and $\rep ak$ means the argument $a$ repeated $k$ times.
It therefore suffices to prove $P_{m,r}\ge0$ on the cube $[-1/2,1/2]^r$.

The positivity of $P_{m,r}$ is proved through three integral representations.
First, up to a positive factor, $h_d(y_1,\ldots,y_k)$ is the $d$-th moment of the convex combination $\sum_i\Theta_iy_i$ whose weight vector $\Theta$ is uniform on the simplex; applying this to the three terms of $P_{m,r}$ at once gives
\[
  P_{m,r}(q;\lambda_1,\ldots,\lambda_r)
  =\E\,\widetilde P_{m,r}\Bigl(q,\sum_{i=1}^r\Theta_i\lambda_i\Bigr),
  \qquad
  \widetilde P_{m,r}(q,\ell)=P_{m,r}(q;\ell,\ldots,\ell).
\]
A convex combination of points of $[-1/2,1/2]$ stays in $[-1/2,1/2]$, so the $r$ variables collapse to one, and it suffices to prove $\widetilde P_{m,r}(q,\ell)\ge0$ for $\ell\in[-1/2,1/2]$.
For $\ell\le0$ this follows from a pointwise bound.
For $\ell>0$, a beta integral converts the diagonal polynomial into
\begin{align*}
  \widetilde P_{m,r}(q,\ell)
  &=\frac{\Gamma(m)}{r\,\Gamma(n)\Gamma(r)^2}\,\ell^{n+r}
  \int_0^\infty\frac{s^{r-1}}{(\ell+s)^m}\,\rho_{n,m}(q+s)\dd s,
  \\
  \rho_{n,m}(u)&=\frac mn\bigl(u^n+(1-u)^n\bigr)-u^{n-1}.
\end{align*}
The function $\rho_{n,m}$ is negative on part of the range of integration, so no pointwise argument applies, and the integral itself has to be bounded from below.
The Laplace representation of $(\ell+s)^{-m}$ replaces the integral by expectations over a gamma random variable $Y\sim\GammaLaw(r,1)$ and reduces everything to
\[
  \E\,\rho_{n,m}(q+xY)\ge0
  \qquad(x\ge0).
\]
Centering at $u=1/2$ and using that $n$ is odd, this becomes a single inequality between consecutive moments of a shifted gamma variable,
\[
  3j\,\E(zY-1)^{2j-1}\le(r+j)\,\E(zY-1)^{2j}
  \qquad(z\ge0),
\]
which we prove from the Stein identity of the gamma distribution and the differential equation it induces.
The constant $3$ is where the hypothesis $p\ge2/3$, that is $q\le1/3$, enters.
The whole argument is uniform in $m$ and involves no computer assistance.
\fi

% TODO(TK): write the organization paragraph once the revised, more
% combinatorial proof and its section structure are in place.
The remainder of the paper proves \Cref{thm:main,thm:commonality}.
\Cref{sec:conclusion} discusses the remaining range $1/2<p<2/3$ and the extension of Goodman-style bounds to other graphs.

\section{Reduction to symmetric matrices}
\label{sec:matrix-reduction}

Let $W\colon[0,1]^2\to[0,1]$ be a graphon of edge density $p\ge2/3$, and
write
\[
  U=1-W,
  \qquad
  q=1-p.
\]
Thus $U$ is the complementary graphon and has edge density $q$.

\begin{lemma}[Step-graphon reduction]
\label{lem:step-reduction}
  It suffices to prove \Cref{thm:main,thm:commonality} for graphons that
  are constant on the rectangles of a partition of $[0,1]$ into equal
  intervals.
\end{lemma}
The proof is given in \Cref{app:step-reduction}.

Fix such a step graphon on $N$ equal intervals.  If $K_{ab}$ is the value
of a step kernel $K$ on $I_a\times I_b$, define the $N\times N$ matrix
$M_K$ and the vector $e\in\R^N$ by
\begin{equation}\label{eq:MK-e}
  M_K=\left(\frac{K_{ab}}N\right)_{a,b=1}^N,
  \qquad
  e=\frac1{\sqrt N}(1,\ldots,1)^{\mathsf T}.
\end{equation}
Then $M_K$ is symmetric and $\norm e=1$.  Also,
\begin{align}
  \Tr(M_K^j)
  &=\frac1{N^j}\sum_{a_1,\ldots,a_j=1}^N
    K_{a_1a_2}\cdots K_{a_ja_1}
   =\dens{C_j}K, \label{eq:cycle-matrix}\\
  e^{\mathsf T}M_K^je
  &=\frac1{N^{j+1}}\sum_{a_0,\ldots,a_j=1}^N
    K_{a_0a_1}\cdots K_{a_{j-1}a_j}
   =\dens{P_j}K, \label{eq:path-matrix}\\
  e^{\mathsf T}M_Ke
  &=\frac1{N^2}\sum_{a,b=1}^NK_{ab}
   =\edens K. \label{eq:edge-matrix}
\end{align}

Let $\mathcal E_0=e^\perp$, and let $I$ be the $N\times N$ identity
matrix.  The matrices
\[
  P=ee^{\mathsf T},
  \qquad
  Q=I-P
\]
are the orthogonal projections of $\R^N$ onto $\R e$ and $\mathcal E_0$,
respectively.  Define the vector $g\in\R^N$ and the symmetric linear map
$A\colon\mathcal E_0\to\mathcal E_0$ by
\begin{equation}\label{eq:g-A}
  g=M_Ue-qe,
  \qquad
  A=QM_UQ\big|_{\mathcal E_0}.
\end{equation}
The edge identity \eqref{eq:edge-matrix} gives
\[
  e^{\mathsf T}g=e^{\mathsf T}M_Ue-q=0,
\]
so $g\in\mathcal E_0$.  Relative to $\R e\oplus\mathcal E_0$,
\begin{equation}\label{eq:block-UW}
  M_U=
  \begin{pmatrix}
    q&g^{\mathsf T}\\
    g&A
  \end{pmatrix},
  \qquad
  M_W=
  \begin{pmatrix}
    p&-g^{\mathsf T}\\
    -g&-A
  \end{pmatrix}.
\end{equation}
The matrix corresponding to the constant graphon $1$ is $M_1=P$, and
\[
  M_W=M_1-M_U.
\]

\begin{lemma}[Eigenvalue bound]\label{lem:half-norm}
  Every eigenvalue of $A$ lies in $[-1/2,1/2]$.
\end{lemma}

\begin{proof}
  Let $f\in\mathcal E_0$ with $\norm f=1$, and let $f_+$ and $f_-$ be its
  coordinatewise positive and negative parts.  Since $e^{\mathsf T}f=0$,
  \[
    e^{\mathsf T}f_+=e^{\mathsf T}f_-=:a,
    \qquad
    2a=e^{\mathsf T}|f|\le\norm e\norm f=1.
  \]
  Expanding $f=f_+-f_-$ gives
  \[
    f^{\mathsf T}M_Uf
    =f_+^{\mathsf T}M_Uf_+
     -f_+^{\mathsf T}M_Uf_-
     -f_-^{\mathsf T}M_Uf_+
     +f_-^{\mathsf T}M_Uf_-.
  \]
  Since $f_+$, $f_-$, and $M_U$ have nonnegative entries,
  $f_+^{\mathsf T}M_Uf_-$ and $f_-^{\mathsf T}M_Uf_+$ are nonnegative.
  Dropping these terms gives
  \[
    f^{\mathsf T}M_Uf
    \le f_+^{\mathsf T}M_Uf_++f_-^{\mathsf T}M_Uf_-
      \le2a^2\le\frac12,
  \]
  where the second inequality follows from $M_U\le M_1$ entrywise and
  $f_+^{\mathsf T}M_1f_+=f_-^{\mathsf T}M_1f_-=a^2$.  Similarly,
  $f_+^{\mathsf T}M_Uf_+$ and $f_-^{\mathsf T}M_Uf_-$ are nonnegative,
  so dropping these terms gives
  \[
    f^{\mathsf T}M_Uf
    \ge-f_+^{\mathsf T}M_Uf_--f_-^{\mathsf T}M_Uf_+
      \ge-2a^2\ge-\frac12.
  \]
  Here the second inequality follows from $M_U\le M_1$ entrywise and
  $f_+^{\mathsf T}M_1f_-=f_-^{\mathsf T}M_1f_+=a^2$.
  Since $f^{\mathsf T}Af=f^{\mathsf T}M_Uf$, the claim follows.
\end{proof}

Let $\lambda_1,\ldots,\lambda_{N-1}$ be the eigenvalues of $A$, counted
with multiplicity, and choose corresponding orthonormal eigenvectors
$u_1,\ldots,u_{N-1}$ of $\mathcal E_0$.  For each $i$, let $P_i$ be the
orthogonal projection onto the one-dimensional subspace spanned by $u_i$.
Define the finite positive measure
\begin{equation}\label{eq:mu}
  \mu=\sum_{i=1}^{N-1}\norm{P_ig}^2\delta_{\lambda_i}.
\end{equation}
The measure $\mu$ is independent of the choice of orthonormal eigenvectors
within each eigenspace.  Here $\delta_{\lambda_i}$ is the unit point mass
at $\lambda_i$.  By \Cref{lem:half-norm},
\begin{equation}\label{eq:moments}
  \supp\mu\subseteq[-1/2,1/2],
  \qquad
  g^{\mathsf T}A^jg=\int\lambda^j\dd\mu(\lambda),
  \qquad
  \mu(\R)=\norm g^2.
\end{equation}

\section{Excursion expansions}
\label{sec:excursions}

For variables $y_1,\ldots,y_k$, the complete homogeneous symmetric
polynomial of degree $d$ is
\begin{equation}\label{eq:h-def}
  h_d(y_1,\ldots,y_k)
  =\sum_{\substack{\gamma_1,\ldots,\gamma_k\ge0\\
                    \gamma_1+\cdots+\gamma_k=d}}
   y_1^{\gamma_1}\cdots y_k^{\gamma_k},
  \qquad h_0=1.
\end{equation}
We write $\rep ak$ when $a$ occurs $k$ times.

\begin{lemma}[Cycle densities by excursions]
\label{lem:cycle-excursions}
  For every odd $m\ge3$,
  \begin{align}
    \tC mU
    ={}&q^m+\Tr(A^m)\notag\\
    &+\sum_{r=1}^{(m-1)/2}\frac mr
      \int\!\cdots\!\int
      h_{m-2r}(\rep qr,\lambda_1,\ldots,\lambda_r)
      \prod_{i=1}^r\dd\mu(\lambda_i),
    \label{eq:cycle-U}\\
    \tC mW
    ={}&p^m-\Tr(A^m)\notag\\
    &+\sum_{r=1}^{(m-1)/2}\frac mr
      \int\!\cdots\!\int
      h_{m-2r}(\rep pr,-\lambda_1,\ldots,-\lambda_r)
      \prod_{i=1}^r\dd\mu(\lambda_i).
    \label{eq:cycle-W}
  \end{align}
\end{lemma}

\begin{proof}
  Since $P+Q=I$,
  \[
    \Tr(M_U^m)
    =\sum_{R_1,\ldots,R_m\in\{P,Q\}}
      \Tr(R_1M_UR_2M_U\cdots R_mM_U).
  \]
  Pair each summand with the word $R_1\cdots R_m$.  If $R_i=P$ for
  every $i$, the summand is $q^m$; if $R_i=Q$ for every $i$, it is
  $\Tr(A^m)$.

  Now consider a mixed word with $r$ blocks of $Q$.  By cyclicity of the
  trace, we may circularly shift the word so that its first letter is $P$
  and its last letter is $Q$.  For $1\le i\le r$, let $\alpha_i+1$ be the
  number of letters in the $i$th $P$-block, and let $\beta_i+1$ be the
  number of letters in the $Q$-block that follows it.  For example, the
  following word has $r=2$:
  \[
    \underbrace{PPP}_{\alpha_1=2}\,
    \underbrace{QQQ}_{\beta_1=2}\,
    \underbrace{PPPP}_{\alpha_2=3}\,
    \underbrace{QQ}_{\beta_2=1}
  \]
  These integers satisfy the length condition, and the value of the
  corresponding summand is determined solely by them:
  \begin{equation}\label{eq:excursion-length-weight}
    \sum_{i=1}^r(\alpha_i+\beta_i)=m-2r,
    \qquad
    \Tr(R_1M_UR_2M_U\cdots R_mM_U)
    =q^{\alpha_1+\cdots+\alpha_r}
     \prod_{i=1}^rg^{\mathsf T}A^{\beta_i}g.
  \end{equation}

  Consider pairs consisting of a word $R_1\cdots R_m$ and an ordered
  tuple $(\alpha_1,\beta_1,\ldots,\alpha_r,\beta_r)$ obtained by
  circularly shifting the word so that it begins with $P$ and ends with
  $Q$, and then taking the successive block lengths.  Equal pairs arising
  from different shifts are counted with multiplicity.  Double counting
  these pairs gives
  \[
    \begin{aligned}
      &\sum_{\substack{\alpha_1,\beta_1,\ldots,\alpha_r,\beta_r\ge0\\
              \sum_{i=1}^r(\alpha_i+\beta_i)=m-2r}}
       \left|\left\{R_1\cdots R_m\in\{P,Q\}^m:
       \substack{(\alpha_1,\beta_1,\ldots,\alpha_r,\beta_r)
       \text{ is obtained}\\
       \text{from a circular shift of }R_1\cdots R_m}
       \right\}\right|\\
      &\qquad=\left|\left\{
       \substack{(R_1\cdots R_m,
       (\alpha_1,\beta_1,\ldots,\alpha_r,\beta_r)):\\
       R_1\cdots R_m\in\{P,Q\}^m,\quad
       \alpha_1,\beta_1,\ldots,\alpha_r,\beta_r\ge0,\\
       \sum_{i=1}^r(\alpha_i+\beta_i)=m-2r,\\
       (\alpha_1,\beta_1,\ldots,\alpha_r,\beta_r)
       \text{ is obtained from a circular shift of }R_1\cdots R_m}
       \right\}\right|\\
      &\qquad=\sum_{\substack{R_1\cdots R_m\in\{P,Q\}^m\\
                       R_1\cdots R_m\text{ has }r\text{ $Q$-blocks}}}
       \left|\left\{\substack{(\alpha_1,\beta_1,\ldots,
       \alpha_r,\beta_r)\in\mathbb Z_{\ge0}^{2r}:\\
       \sum_{i=1}^r(\alpha_i+\beta_i)=m-2r,\\
       (\alpha_1,\beta_1,\ldots,\alpha_r,\beta_r)
       \text{ is obtained from a circular shift of }R_1\cdots R_m}
       \right\}\right|.
    \end{aligned}
  \]
  For every ordered tuple in the first sum,
  \[
    \left|\left\{R_1\cdots R_m\in\{P,Q\}^m:
      \substack{(\alpha_1,\beta_1,\ldots,\alpha_r,\beta_r)
      \text{ is obtained}\\
      \text{from a circular shift of }R_1\cdots R_m}
    \right\}\right|=m.
  \]
  For every word in the last sum,
  \[
    \left|\left\{\substack{(\alpha_1,\beta_1,\ldots,
      \alpha_r,\beta_r)\in\mathbb Z_{\ge0}^{2r}:\\
      \sum_{i=1}^r(\alpha_i+\beta_i)=m-2r,\\
      (\alpha_1,\beta_1,\ldots,\alpha_r,\beta_r)
      \text{ is obtained from a circular shift of }R_1\cdots R_m}
    \right\}\right|=r.
  \]
  Applying the same double count with the values in
  \eqref{eq:excursion-length-weight} as weights gives
  \begin{equation}\label{eq:U-r-excursions}
    \begin{aligned}
      &\sum_{\substack{R_1\cdots R_m\in\{P,Q\}^m\\
                       R_1\cdots R_m\text{ has }r\text{ $Q$-blocks}}}
        \Tr(R_1M_U\cdots R_mM_U)\\
      &\qquad=\frac mr
        \sum_{\substack{\alpha_1,\beta_1,\ldots,
                          \alpha_r,\beta_r\ge0\\
                          \sum_{i=1}^r(\alpha_i+\beta_i)=m-2r}}
        q^{\alpha_1+\cdots+\alpha_r}
        \prod_{i=1}^rg^{\mathsf T}A^{\beta_i}g.
    \end{aligned}
  \end{equation}
  Summing \eqref{eq:U-r-excursions} over $r$ and including the two
  constant words gives
  \[
    \tC mU=q^m+\Tr(A^m)
    +\sum_{r=1}^{(m-1)/2}\frac mr
      \sum_{\substack{\alpha_1,\beta_1,\ldots,
                        \alpha_r,\beta_r\ge0\\
                        \sum_{i=1}^r(\alpha_i+\beta_i)=m-2r}}
      q^{\alpha_1+\cdots+\alpha_r}
      \prod_{i=1}^rg^{\mathsf T}A^{\beta_i}g.
  \]
  Applying the moment identity \eqref{eq:moments} now gives
  \[
    \begin{aligned}
      \tC mU={}&q^m+\Tr(A^m)\\
      &+\sum_{r=1}^{(m-1)/2}\frac mr
      \int\!\cdots\!\int
      \sum_{\substack{\alpha_1,\beta_1,\ldots,
                        \alpha_r,\beta_r\ge0\\
                        \sum_{i=1}^r(\alpha_i+\beta_i)=m-2r}}
      q^{\alpha_1+\cdots+\alpha_r}
      \lambda_1^{\beta_1}\cdots\lambda_r^{\beta_r}
      \prod_{i=1}^r\dd\mu(\lambda_i).
    \end{aligned}
  \]
  The inner sum is
  $h_{m-2r}(\rep qr,\lambda_1,\ldots,\lambda_r)$ by
  \eqref{eq:h-def}.  This proves \eqref{eq:cycle-U}.

  For $M_W$, the block identity \eqref{eq:block-UW} replaces
  $(q,g,A)$ by $(p,-g,-A)$.  The two crossing signs cancel in each block,
  while
  \[
    (-g)^{\mathsf T}(-A)^{\beta_i}(-g)
    =(-1)^{\beta_i}g^{\mathsf T}A^{\beta_i}g.
  \]
  Hence $q$ becomes $p$, every $\lambda_i$ becomes $-\lambda_i$, and
  $\Tr(A^m)$ becomes $-\Tr(A^m)$.  This proves \eqref{eq:cycle-W}.
\end{proof}

Let $P_k$ denote the path with $k$ edges, and define the scalar
\begin{equation}\label{eq:xk}
  x_k=\tP kU=e^{\mathsf T}M_U^ke.
\end{equation}

\begin{lemma}[Path densities by excursions]
\label{lem:path-excursions}
  For every $k\ge0$,
  \begin{equation}\label{eq:path-expansion}
    x_k=q^k+
    \sum_{r=1}^{\lfloor k/2\rfloor}
      \int\!\cdots\!\int
      h_{k-2r}(\rep q{r+1},\lambda_1,\ldots,\lambda_r)
      \prod_{i=1}^r\dd\mu(\lambda_i).
  \end{equation}
\end{lemma}

\begin{proof}
  The case $k=0$ is immediate.  Since $Pe=e$ and $Qe=0$, a word that
  starts or ends with $Q$ has summand zero.  Hence, for $k\ge1$, inserting
  $P+Q=I$ between consecutive factors of $M_U$ gives
  \[
    x_k=
    \sum_{(R_1,\ldots,R_{k-1})\in\{P,Q\}^{k-1}}
      e^{\mathsf T}P M_U R_1M_U\cdots R_{k-1}M_UPe.
  \]
  Pair each summand with the word $PR_1\cdots R_{k-1}P$.  If every letter
  is $P$, the summand is $q^k$.

  Now consider a word with $r$ blocks of $Q$.  Let
  $\alpha_0+1,\ldots,\alpha_r+1$ be the numbers of letters in the
  successive $P$-blocks, and let $\beta_1+1,\ldots,\beta_r+1$ be the
  numbers of letters in the successive $Q$-blocks.  For example, the
  following word has $r=2$:
  \[
    \underbrace{PP}_{\alpha_0=1}\,
    \underbrace{QQ}_{\beta_1=1}\,
    \underbrace{PP}_{\alpha_1=1}\,
    \underbrace{Q}_{\beta_2=0}\,
    \underbrace{PP}_{\alpha_2=1}.
  \]
  These integers satisfy the length condition, and the corresponding
  summand is
  \[
    \sum_{i=0}^r\alpha_i+\sum_{i=1}^r\beta_i=k-2r,
    \qquad
    e^{\mathsf T}P M_U R_1M_U\cdots R_{k-1}M_UPe
    =q^{\alpha_0+\cdots+\alpha_r}
      \prod_{i=1}^rg^{\mathsf T}A^{\beta_i}g.
  \]
  There is a one-to-one correspondence between such words and the values
  of the $\alpha_i$ and $\beta_i$ satisfying the length condition.
  Therefore,
  \begin{equation}\label{eq:path-r-excursions}
    \begin{aligned}
      &\sum_{\substack{(R_1,\ldots,R_{k-1})\in\{P,Q\}^{k-1}\\
                       PR_1\cdots R_{k-1}P
                       \text{ has }r\text{ $Q$-blocks}}}
        e^{\mathsf T}P M_U R_1M_U\cdots R_{k-1}M_UPe\\
      &\qquad=
      \sum_{\substack{\alpha_0,\ldots,\alpha_r,\beta_1,\ldots,\beta_r\ge0\\
                      \sum_{i=0}^r\alpha_i+
                      \sum_{i=1}^r\beta_i=k-2r}}
      q^{\alpha_0+\cdots+\alpha_r}
      \prod_{i=1}^rg^{\mathsf T}A^{\beta_i}g.
    \end{aligned}
  \end{equation}
  Summing \eqref{eq:path-r-excursions} over $r$ and including the word
  consisting only of $P$ gives
  \[
    x_k=q^k+
    \sum_{r=1}^{\lfloor k/2\rfloor}
      \sum_{\substack{\alpha_0,\ldots,\alpha_r,\beta_1,\ldots,\beta_r\ge0\\
                      \sum_{i=0}^r\alpha_i+
                      \sum_{i=1}^r\beta_i=k-2r}}
      q^{\alpha_0+\cdots+\alpha_r}
      \prod_{i=1}^rg^{\mathsf T}A^{\beta_i}g.
  \]
  Applying the moment identity \eqref{eq:moments} now gives
  \[
    \begin{aligned}
      x_k={}&q^k+
      \sum_{r=1}^{\lfloor k/2\rfloor}
      \int\!\cdots\!\int
      \sum_{\substack{\alpha_0,\ldots,\alpha_r,\beta_1,\ldots,\beta_r\ge0\\
                      \sum_{i=0}^r\alpha_i+
                      \sum_{i=1}^r\beta_i=k-2r}}
      q^{\alpha_0+\cdots+\alpha_r}
      \lambda_1^{\beta_1}\cdots\lambda_r^{\beta_r}
      \prod_{i=1}^r\dd\mu(\lambda_i).
    \end{aligned}
  \]
  The inner sum is
  $h_{k-2r}(\rep q{r+1},\lambda_1,\ldots,\lambda_r)$ by
  \eqref{eq:h-def}.  This proves \eqref{eq:path-expansion}.
\end{proof}

\subsection{Reduction to symmetric polynomials}

\begin{lemma}[Reduction to symmetric polynomials]
\label{lem:polynomial-reduction}
  For $1\le r\le(m-1)/2$, define the symmetric polynomial
  \begin{align}
    P_{m,r}(q;\lambda_1,\ldots,\lambda_r)
    ={}&\frac mr\Bigl[
      h_{m-2r}(\rep pr,-\lambda_1,\ldots,-\lambda_r)
      +h_{m-2r}(\rep qr,\lambda_1,\ldots,\lambda_r)
    \Bigr]\notag\\
    &-h_{m-2r-1}(\rep q{r+1},\lambda_1,\ldots,\lambda_r).
    \label{eq:Pmr}
  \end{align}
  It suffices to prove
  \begin{equation}\label{eq:Phi-expansion}
    \Phi_m:=\sum_{r=1}^{(m-1)/2}
      \int\!\cdots\!\int
      P_{m,r}(q;\lambda_1,\ldots,\lambda_r)
      \prod_{i=1}^r\dd\mu(\lambda_i)\ge0
  \end{equation}
  to prove \Cref{thm:main,thm:commonality}.
\end{lemma}

\begin{proof}
  Adding the cycle expansions \eqref{eq:cycle-U} and
  \eqref{eq:cycle-W} cancels the two trace terms and gives
  \[
    \begin{aligned}
      \tC mW&+\tC mU\\
      =p^m+q^m&+\sum_{r=1}^{(m-1)/2}\frac mr
      \int\!\cdots\!\int
      [
        h_{m-2r}(\rep pr,-\lambda_1,\ldots,-\lambda_r)
        +h_{m-2r}(\rep qr,\lambda_1,\ldots,\lambda_r)
      ]
      \prod_{i=1}^r\dd\mu(\lambda_i).
    \end{aligned}
  \]
  The path expansion \eqref{eq:path-expansion} with $k=m-1$ is
  \[
    x_{m-1}=q^{m-1}+
    \sum_{r=1}^{(m-1)/2}
      \int\!\cdots\!\int
      h_{m-2r-1}(\rep q{r+1},\lambda_1,\ldots,\lambda_r)
      \prod_{i=1}^r\dd\mu(\lambda_i).
  \]
  For each $r$, subtracting the integrand in the path expansion from the
  corresponding integrand in the sum of the cycle expansions leaves
  \[
    \frac mr\Bigl[
      h_{m-2r}(\rep pr,-\lambda_1,\ldots,-\lambda_r)
      +h_{m-2r}(\rep qr,\lambda_1,\ldots,\lambda_r)
    \Bigr]
    -h_{m-2r-1}(\rep q{r+1},\lambda_1,\ldots,\lambda_r),
  \]
  which is $P_{m,r}(q;\lambda_1,\ldots,\lambda_r)$ by
  \eqref{eq:Pmr}.  Therefore,
  \begin{equation}\label{eq:commonality-gap}
    \tC mW+\tC mU-x_{m-1}
    =p^m+q^m-q^{m-1}+\Phi_m.
  \end{equation}
  Hence $\Phi_m\ge0$ gives
  \[
    \tC mW+\tC mU
    \ge p^m+q^m+\tP{m-1}U-q^{m-1},
  \]
  proving \Cref{thm:commonality}.

  Since $0\le U\le1$, deleting the edge joining $x_m$ and $x_1$ gives
  \begin{equation}\label{eq:cycle-path}
    \begin{aligned}
      \tC mU
      &=\int_{[0,1]^m}
        U(x_1,x_2)\cdots U(x_{m-1},x_m)U(x_m,x_1)
        \dd x_1\cdots\dd x_m\\
      &\le\int_{[0,1]^m}
        U(x_1,x_2)\cdots U(x_{m-1},x_m)
        \dd x_1\cdots\dd x_m
       =x_{m-1}.
    \end{aligned}
  \end{equation}
  Rearranging \eqref{eq:commonality-gap} and using $p+q=1$ gives
  \begin{align*}
    \tC mW-\bigl(p^m-pq^{m-1}\bigr)
    &={}
      \Phi_m+x_{m-1}-\tC mU
      +q^m-q^{m-1}+pq^{m-1}\\
    &=\Phi_m+x_{m-1}-\tC mU
     \ge\Phi_m,
  \end{align*}
  where the last inequality follows from \eqref{eq:cycle-path}.
  Thus $\Phi_m\ge0$ also proves \Cref{thm:main}.
\end{proof}

Since $\mu$ is positive and supported on $[-1/2,1/2]$, it suffices to
prove
\[
  P_{m,r}(q;\lambda_1,\ldots,\lambda_r)\ge0
\]
for $0\le q\le1/3$, $1\le r\le(m-1)/2$, and
$\lambda_i\in[-1/2,1/2]$.

\section{Conclusion} \label{sec:conclusion}

The remaining question is the intermediate range $p \in (1/2, 2/3)$, which was not covered by our result.
The proof technique used in the range $p \ge 2/3$ is not applicable to the intermediate range: the edge-deletion lemma is too loose in this range and makes the inequality invalid.
We have proven the same lower bound on the range $p \in (1/2, 2/3)$, but the proof is not combinatorial and we are in the process of polishing it to make it more readable.
We note that the stronger extended commonality bound is not valid in the range $p \in (1/2, 2/3)$, as the edge deletion lemma gap is exactly the same as the gap the extended commonality bound requires.

The definition of the Goodman-style bound can be extended to other graphs.
For a graph $H$, a Goodman-style bound is the following lower bound on the density of $H$ at the edge densities $p \ge 1 - \frac{1}{\chi(H) - 1}$:
\[
  t(H, W) \ge \chi_H\left( \frac{1}{1-p} \right) (1-p)^{v(H)},
\]
where $\chi_H$ is the chromatic polynomial of $H$, $\chi(H)$ is the chromatic number of $H$, and $v(H)$ is the number of vertices in $H$.
If the edge-density-constrained minimisation problem
\[
  \min_W t(H, W) \text{ s.t. } t(K_2, W) = p
\]
is minimised by the balanced complete $k$-partite graph at $p = 1 - 1 / k$ for all $k \ge \chi(H) - 1$, then the Goodman-style bound is tight at these points, and is the unique polynomial that interpolates these points.

Thus a general research question is determining which graphs have the Goodman-style bound and which graphs do not have the Goodman-style bound.
We list some conditions that partially characterise the nonbipartite graphs with the Goodman-style bound.
\begin{itemize}
  \item If a graph is chordal and all its maximal cliques are of the same size, then it has the Goodman-style bound.
  \item If a graph has an even girth, then it does not have the Goodman-style bound.
\end{itemize}

\section{Remarks} \label{sec:remarks}

\textbf{AI Disclosure} 
This proof was discovered by two week span of interaction with ChatGPT 5.5 and Fable 5. 
We polished the proof and modified the proof to give combinatorial interpretation and make it more readable.

\tk{TODO: Write here about deviation that lean verification and current proof.}

\appendix

\section{Step-graphon reduction}
\label{app:step-reduction}

\begin{proof}[Proof of \Cref{lem:step-reduction}]
  For $N=2^n$, let $I_1,\ldots,I_N$ be the intervals of length $1/N$, and
  define
  \[
    W_n(x,y)=N^2\int_{I_a\times I_b}W(s,t)\dd s\dd t
    \qquad ((x,y)\in I_a\times I_b).
  \]
  Each $W_n$ is a symmetric graphon, and
  \[
    \int_{[0,1]^2}W_n
    =\sum_{a,b=1}^N\int_{I_a\times I_b}W=p.
  \]
  Moreover,
  \begin{equation}\label{eq:step-L1}
    \norm{W_n-W}_1\longrightarrow0.
  \end{equation}
  To verify the convergence, let $S$ be a function constant on the
  level-$n_0$ dyadic rectangles.  For $n\ge n_0$, cell averaging fixes
  $S$ and is an $L^1$-contraction, so
  \[
    \norm{W_n-W}_1
    \le \norm{W_n-S}_1+\norm{S-W}_1
    \le2\norm{W-S}_1.
  \]
  Dyadic step functions are dense in $L^1([0,1]^2)$, proving
  \eqref{eq:step-L1}.

  If $F$ is a finite graph, telescoping its edge factors gives
  \[
    \left|\dens F{W_n}-\dens FW\right|
    \le |E(F)|\norm{W_n-W}_1.
  \]
  The same bound holds for $U_n=1-W_n$ and $U=1-W$.  Hence all cycle and
  path densities in the two inequalities converge, while $p$ and $q$
  remain unchanged.
\end{proof}

\bibliographystyle{abbrv}
\bibliography{references}

\end{document}
